\input{configTD}

\title{Analyse Numérique TD3 - Méthodes itératives et matrices stochastiques}

\author{}
\date{}

\newenvironment{solution}
    {\par\vspace{0.5em}\begin{mdframed}[linewidth=0.5pt]\par}
    {\end{mdframed}\par\vspace{0.5em}}

\begin{document}
\sffamily
\maketitle

\section*{Exercice 1 - Méthode de Jacobi}
On considère le système
\[
\begin{cases}
4x_1-x_2=7,\\
-x_1+4x_2-x_3=6,\\
-x_2+4x_3=7.
\end{cases}
\]

\begin{enumerate}
    \item Écrire ce système sous la forme $Ax=b$.

    \item Écrire les formules de Jacobi permettant de construire $x^{(k+1)}$ à partir de $x^{(k)}$.

    \item À partir de $x^{(0)}=(0,0,0)^T$, calculer $x^{(1)}$ puis $x^{(2)}$.

    \item Expliquer pourquoi on peut raisonnablement s'attendre à la convergence.

    \ifthenelse{\boolean{showSolutions}}{
        \begin{solution}
            La matrice est
            \[
            A=\begin{pmatrix}
            4 & -1 & 0\\
            -1 & 4 & -1\\
            0 & -1 & 4
            \end{pmatrix},
            \qquad
            b=\begin{pmatrix}
            7\\6\\7
            \end{pmatrix}.
            \]
            Les équations de Jacobi sont
            \[
            x_1^{(k+1)}=\frac{7+x_2^{(k)}}{4},
            \qquad
            x_2^{(k+1)}=\frac{6+x_1^{(k)}+x_3^{(k)}}{4},
            \qquad
            x_3^{(k+1)}=\frac{7+x_2^{(k)}}{4}.
            \]
            Avec $x^{(0)}=(0,0,0)^T$,
            \[
            x^{(1)}=\begin{pmatrix}
            7/4\\
            3/2\\
            7/4
            \end{pmatrix}.
            \]
            Puis
            \[
            x^{(2)}=\begin{pmatrix}
            17/8\\
            19/8\\
            17/8
            \end{pmatrix}.
            \]
            La matrice est à diagonale strictement dominante, donc la méthode de Jacobi converge.
        \end{solution}
    }{}
\end{enumerate}

\section*{Exercice 2 - Jacobi ou divergence ?}
On considère maintenant le système
\[
\begin{cases}
x_1+3x_2=4,\\
3x_1+x_2=4.
\end{cases}
\]

\begin{enumerate}
    \item Écrire l'itération de Jacobi.

    \item À partir de $x^{(0)}=(0,0)^T$, calculer $x^{(1)}$, $x^{(2)}$ et $x^{(3)}$.

    \item Que constate-t-on ?

    \item Calculer la solution exacte et comparer avec le comportement itératif.

    \ifthenelse{\boolean{showSolutions}}{
        \begin{solution}
            L'itération de Jacobi est
            \[
            x_1^{(k+1)}=4-3x_2^{(k)},
            \qquad
            x_2^{(k+1)}=4-3x_1^{(k)}.
            \]
            À partir de $(0,0)$ :
            \[
            x^{(1)}=(4,4),\qquad
            x^{(2)}=(-8,-8),\qquad
            x^{(3)}=(28,28).
            \]
            La suite diverge.
            Pourtant la solution exacte existe :
            \[
            x_1=x_2=1.
            \]
            Le problème ne vient donc pas de l'existence de la solution, mais du choix de la méthode itérative sur ce système.
        \end{solution}
    }{}
\end{enumerate}

\section*{Exercice 3 - PageRank simplifié}
On considère trois pages web $A$, $B$ et $C$.
Le lien sortant de $A$ pointe vers $B$ et $C$ avec probabilité $1/2$ chacune.
Le lien sortant de $B$ pointe toujours vers $C$.
Le lien sortant de $C$ pointe toujours vers $A$.

On note
\[
P=
\begin{pmatrix}
0 & 0 & 1\\
1/2 & 0 & 0\\
1/2 & 1 & 0
\end{pmatrix}.
\]
La distribution stationnaire $\pi=(\pi_A,\pi_B,\pi_C)^T$ vérifie
\[
P\pi=\pi,
\qquad
\pi_A+\pi_B+\pi_C=1.
\]

\begin{enumerate}
    \item Écrire le système vérifié par $\pi$.

    \item Déterminer la distribution stationnaire.

    \item Quelle page obtient le score le plus élevé ?

    \ifthenelse{\boolean{showSolutions}}{
        \begin{solution}
            Le système est
            \[
            \pi_A=\pi_C,\qquad
            \pi_B=\frac12\pi_A,\qquad
            \pi_C=\frac12\pi_A+\pi_B.
            \]
            Avec la condition de normalisation,
            \[
            \pi_A+\frac12\pi_A+\pi_A=1
            \]
            donc
            \[
            \frac52\pi_A=1,\qquad \pi_A=\frac25.
            \]
            Ainsi
            \[
            \pi_B=\frac15,\qquad \pi_C=\frac25.
            \]
            Les pages $A$ et $C$ obtiennent ici le même score maximal.
        \end{solution}
    }{}
\end{enumerate}

\section*{Exercice 4 - Matrice stochastique et classes de risque}
Un assureur classe un portefeuille en trois catégories de risque : faible ($F$), moyen ($M$), élevé ($E$).
D'une année à l'autre, la probabilité de changer de classe est modélisée par
\[
T=
\begin{pmatrix}
0.8 & 0.2 & 0\\
0.1 & 0.7 & 0.2\\
0 & 0.3 & 0.7
\end{pmatrix}.
\]
Le vecteur ligne $p_n=(f_n,m_n,e_n)$ représente la répartition l'année $n$.

\begin{enumerate}
    \item Vérifier que $T$ est stochastique.

    \item Si
    \[
    p_0=(0.6,0.3,0.1),
    \]
    calculer $p_1$ puis $p_2$.

    \item On associe les primes annuelles suivantes :
    \[
    800\ \text{euros},\qquad 1200\ \text{euros},\qquad 1800\ \text{euros}.
    \]
    Donner la prime moyenne attendue après une année.

    \item Expliquer en quelques lignes l'intérêt d'un tel modèle.

    \ifthenelse{\boolean{showSolutions}}{
        \begin{solution}
            Chaque ligne somme à $1$, donc $T$ est bien stochastique.
            On calcule
            \[
            p_1=p_0T=(0.51,0.40,0.09).
            \]
            Puis
            \[
            p_2=p_1T=(0.448,0.409,0.143).
            \]
            La prime moyenne après une année vaut
            \[
            0.51\times 800+0.40\times 1200+0.09\times 1800=1050\ \text{euros}.
            \]
            Ce modèle permet d'anticiper l'évolution du portefeuille, de chiffrer une prime moyenne attendue et d'étudier un régime stationnaire à long terme.
        \end{solution}
    }{}
\end{enumerate}

\section*{Exercice 5 - Application à une discrétisation}
On retrouve la matrice
\[
A=\begin{pmatrix}
2 & -1 & 0 & 0\\
-1 & 2 & -1 & 0\\
0 & -1 & 2 & -1\\
0 & 0 & -1 & 2
\end{pmatrix}.
\]

\begin{enumerate}
    \item Montrer que $A$ est symétrique définie positive.

    \item Expliquer pourquoi cette propriété est favorable à l'utilisation de méthodes itératives.

    \item Citer un contexte de discrétisation dans lequel une telle matrice apparaît naturellement.

    \ifthenelse{\boolean{showSolutions}}{
        \begin{solution}
            La matrice est symétrique. De plus, pour tout vecteur $x=(x_1,\dots,x_4)^T$,
            \[
            x^TAx=x_1^2+(x_1-x_2)^2+(x_2-x_3)^2+(x_3-x_4)^2+x_4^2>0
            \]
            pour tout $x\neq 0$.
            Donc $A$ est définie positive.
            Cette structure favorise la convergence de plusieurs méthodes itératives et garantit une solution unique.
            Une telle matrice apparaît naturellement dans les discrétisations 1D de type chaleur, diffusion ou barre élastique.
        \end{solution}
    }{}
\end{enumerate}

\end{document}
